\section{SEC-007: SSRF via \texttt{curl}/\texttt{wget} in CLI Tool}
\label{sec:sec-007}

\subsection*{Metadata}
\begin{tabular}{ll}
\textbf{Issue ID} & SEC-007 \\
\textbf{Severity} & Medium (CVSS 7.5) \\
\textbf{Category} & Security \\
\textbf{CWE} & CWE-918: Server-Side Request Forgery (SSRF) \\
\textbf{Affected Files} & \srcfile{src/tools/shared.ts}, \srcfile{src/tools/cli.ts}, \srcfile{src/tools/sandbox-cli.ts} \\
\textbf{Branch} & \branchref{fix/SEC-007-ssrf-prevention} \\
\textbf{Changes} & \branchdiff{fix/SEC-007-ssrf-prevention} \\
\textbf{Commit} & 6142a53 \\
\end{tabular}

\subsection*{Executive Summary}
The \texttt{cli} tool gives the LLM full shell access via \texttt{spawn("sh", ["-c", command])}. While this enables arbitrary command execution (SEC-001), a specific high-impact subclass is Server-Side Request Forgery (SSRF)---the ability to make the agent process issue HTTP requests to internal network services. When the LLM is instructed to run commands like \texttt{curl http://169.254.169.254/latest/meta-data/}, the agent becomes a pivot point for network reconnaissance and attack.

The fix adds a \texttt{checkSSRF()} function that inspects commands for network tool usage and blocks requests to private IP ranges, internal hostnames, known metadata endpoints, the Docker socket via \texttt{--unix-socket}, and private IPv6 addresses.

\subsection*{Root Cause Analysis}
The \texttt{spawn("sh", ["-c", command])} pattern allows the LLM to run any shell command, including network requests. The tool had no URL allowlist, IP address validation, internal network detection, or protocol restriction. High-value SSRF targets include cloud metadata services (AWS \texttt{169.254.169.254}, GCP \texttt{metadata.google.internal}), Kubernetes API servers (\texttt{kubernetes.default.svc}), internal microservices, and the Docker socket.

\subsection*{Fix Implementation}
The fix adds a \texttt{checkSSRF()} function in \srcfile{src/tools/shared.ts} that:

\begin{enumerate}
    \item Detects the use of network tools (\texttt{curl}, \texttt{wget}, \texttt{fetch}) via \texttt{NETWORK\_TOOLS\_PATTERN}.
    \item Extracts all URLs from the command using a regex that handles IPv6 bracket notation and \texttt{user:pass@host} authentication.
    \item Checks each URL's hostname against private IPv4 ranges (\texttt{127.x.x.x}, \texttt{10.x.x.x}, \texttt{172.16-31.x.x}, \texttt{192.168.x.x}, \texttt{169.254.x.x}, \texttt{0.x.x.x}, \texttt{100.64-127.x.x}).
    \item Checks each hostname against blocked patterns (cloud metadata endpoints, Kubernetes internal DNS, \texttt{.internal} domains, \texttt{localhost}).
    \item Checks for private IPv6 addresses (loopback \texttt{::1}, unique-local \texttt{fc00::/7}, link-local \texttt{fe80::}).
    \item Blocks Docker socket access via \texttt{curl --unix-socket /var/run/docker.sock}.
\end{enumerate}

The function is called at the start of both \texttt{cli.ts} and \texttt{sandbox-cli.ts} execute methods:

\begin{lstlisting}[language=TypeScript]
try {
    checkSSRF(command);
} catch (err) {
    reject(err);
    return;
}
\end{lstlisting}

\subsection*{Affected Files}
\begin{itemize}
    \item \srcfile{src/tools/shared.ts} --- Added \texttt{checkSSRF()} with private IP range patterns, blocked hostname patterns, network tool detection, URL extraction (IPv6 brackets and auth-supporting), Docker socket detection, and IPv6 private address detection.
    \item \srcfile{src/tools/cli.ts} --- Added \texttt{checkSSRF()} call before spawning shell commands.
    \item \srcfile{src/tools/sandbox-cli.ts} --- Added \texttt{checkSSRF()} call before sandbox shell commands.
\end{itemize}

\subsection*{Verification}
The test suite in \texttt{test/ssrf.test.ts} covers 24 test cases:

\begin{itemize}
    \item \textbf{Private IPv4 (7 tests)}: 127.x.x.x, 10.x.x.x, 172.16-31.x.x, 192.168.x.x, 169.254.x.x (link-local), 0.x.x.x, 100.64-127.x.x (CGNAT).
    \item \textbf{Hostnames (5 tests)}: metadata.google.internal, metadata.google.compute, kubernetes.default.svc, .internal domains, localhost.
    \item \textbf{IPv6 (3 tests)}: Loopback [::1], unique-local fc00::/7 and fd00::/7, link-local fe80::.
    \item \textbf{Docker socket (1 test)}: Blocking curl --unix-socket to docker.sock.
    \item \textbf{Allowed (4 tests)}: Public IP, public domain, custom port, subdomain.
    \item \textbf{Non-network (1 test)}: Commands without curl/wget/fetch pass through.
    \item \textbf{Edge cases (3 tests)}: curl -L to private IP, wget with auth in URL (user:pass@host), HTTPS to private IP.
\end{itemize}

\subsection*{Test File}
\srcfile{test/ssrf.test.ts}
